\TieudeT{PHONG TỎA VẬT LÍ}
\subsubsection*{PHẦN I. Thí sinh trả lời câu hỏi từ câu 1 đến câu 18. Mỗi câu hỏi thí sinh chỉ chọn một phương án}
\setcounter{ex}{0}
%mac
\loai{Định luật 3}
\begin{ex}%book
Phát biểu nào sau đây là đúng về định luật III Newton?

\choice
{Khi vật $A$ tác dụng lên vật $B$ một lực, thì vật $B$ cũng tác dụng trở lại vật $A$ một lực. Hai lực này là hai lực trực đối: $\vec{F}_{AB}=\vec{F}_{BA}$}
{Khi vật $A$ tác dụng lên vật $B$ một lực, thì vật $B$ không tác dụng trở lại vật $A$ một lực}
{Khi vật $A$ tác dụng lên vật $B$ một lực, thì vật $B$ cũng tác dụng trở lại vật $A$ một lực. Hai lực này là hai lực cân bằng nhau: $\vec{F}_{AB}=\vec{F}_{BA}$}
{\True Khi vật $A$ tác dụng lên vật $B$ một lực, thì vật $B$ cũng tác dụng trở lại vật $A$ một lực. Hai lực này là hai lực trực đối: $\vec{F}_{AB}+\vec{F}_{BA}=\vec 0$}
\loigiai{
\begin{itemize}
\item Định luật 3 - Newton: Khi vật A tác dụng lên vật $B$ một lực, thì vật $B$ cũng tác dụng trở lại vật $A$ một lực.
\item 	Hai lực này là hai lực trực đối: $\vec{F}_{AB}=-\vec{F}_{BA}$
\item 	Hai lực trực đối cân bằng là hai lực có cùng độ lớn, cùng phương, ngược chiều tác dụng vào hai vật khác nhau
\end{itemize}

}
\haidong{8} \end{ex}
\loai{Lực và phản lực}
\begin{ex}%book
Theo định luật 3 Newton thì lực và phản lực là cặp lực
\choice
{cân bằng}
{có cùng điểm đặt}
{cùng phương, cùng chiều và cùng độ lớn}
{\True xuất hiện và mất đi đồng thời}
\loigiai{
Chọn D

}
\haidong{8} \end{ex}
%\newpage
\begin{ex}%book
Theo định luật $3$ Newton, cặp "lực và phản lực"
\choice
{tác dụng vào cùng một vật}
{\True tác dụng vào hai vật khác nhau}
{không bằng nhau về độ lớn}
{bằng nhau về độ lớn nhưng không cùng giá}
\loigiai{
Chọn B

}
\haidong{8} \end{ex}

\begin{ex}%book
Cặp lực và phản lực trong định luật $3$ Newton
\choice
{không cùng bản chất}
{\True cùng bản chất}
{tác dụng vào cùng một vật}
{bằng nhau về độ lớn nhưng không cùng giá}
\loigiai{
Chọn B

}
\haidong{8} \end{ex}
\begin{ex}%book
Phát biểu nào sau đây là \textbf{sai}? Lực và phản lực
\choice
{\True là hai lực cân bằng}
{luôn xuất hiện đồng thời}
{cùng phương}
{cùng bản chất}
\loigiai{
Chọn A

Lực và phản lực chỉ là hai lực trực đối $\rightarrow$ A \textbf{sai}.

}
\haidong{8} \end{ex}
\begin{ex}%book
Phát biểu nào sau đây là \textbf{sai}? Lực và phản lực
\choice
{là hai lực trực đối}
{cùng độ lớn}
{ngược chiều nhau}
{\True có thể tác dụng vào cùng một vật}
\loigiai{
Lực và phản lực tác dụng vào hai vật khác nhau $\rightarrow D$ \textbf{sai}.

}
\haidong{8} \end{ex}

\begin{ex}%book
Cặp lực và phản lực \textbf{không} phải là hai lực cân bằng vì
\choice
{\True điểm đặt của chúng ở trên hai vật khác nhau}
{điểm đặt của chúng ở trên hai vật giống nhau nhau}
{chúng có độ lớn không bằng nhau}
{chúng có cùng hướng và có độ lớn khác nhau}
\loigiai{
Chọn A

Lực và phản lực không cân bằng nhau (vì chúng đặt vào hai vật khác nhau).

}
\haidong{8} %<MyLT1>
\end{ex}
\loai{Vận dụng định luật 3}
\begin{ex}%book%[1P2Y4-1][Loai1]
Khi dùng búa đóng một cây đinh vào một khúc gỗ thì
\choice
{lực của búa tác dụng vào định lớn hơn lực đinh tác dụng vào búa}
{tùy thuộc đinh di chuyển nhiều hay ít mà lực do đinh tác dụng vào búa lớn hơn hay nhỏ hơn lực do búa tác dụng vào đinh}
{lực của búa tác dụng vào đinh nhỏ hơn lực đinh tác dụng vào búa}
{\True lực của búa tác dụng vào đinh về độ lớn bằng lực của đinh tác dụng vào búa}
\loigiai{
Lực do búa tác dụng vào đinh và lực do đinh tác dụng vào búa là hai lực trực đối $\Rightarrow$ lực của búa tác dụng vào đinh về độ lớn bằng lực của đinh tác dụng vào búa.
}
\haidong{8} \end{ex}

\begin{ex}%book
Trong một cơn giông, một cành cây bị gãy và bay trúng vào một cửa kính, làm vỡ kính thì
\choice
{lực của cành cây tác dụng lên tấm kính lớn hơn lực của tấm kính tác dụng vào cành cây}
{\True lực của cành cây tác dụng lên tấm kính có độ lớn bằng lực của tấm kính tác dụng vào cành cây}
{lực của cành cây tác dụng lên tấm kính nhỏ hơn lực của tấm kính tác dụng vào cành cây}
{cành cây không tương tác với tấm kính khi làm vỡ kính}
\loigiai{
\begin{itemize}
\item Định luật 3 Newton
\item 	Trong mọi trường hợp, khi vật $A$ tác dụng lên vật $B$ một lực thì vật $B$ cũng tác dụng trở lại vật $A$ một lực. Hai lực này là hai lực trực đối
\item 	$\vec{F}_{AB}=-\vec{F}_{BA}$
\item 	Hai lực trực đối là hai lực tác dụng theo cùng một đường thẳng, ngược chiều nhau, có độ lớn bằng nhau và điểm đặt lên hai vật khác nhau
\end{itemize}
}
\haidong{8} \end{ex}
\begin{ex}%book
Một người thực hiện động tác nằm sấp, chống tay xuống sàn nhà để nâng người lên. Hỏi sàn nhà đẩy người đó như thế nào?
\choice
{Không đẩy gì cả}
{Đẩy xuống}
{\True Đẩy lên}
{Đẩy sang bên}
\loigiai{
}
\haidong{8} \end{ex}
\begin{ex}%book
Người ta dùng búa đóng một cây đinh vào một khối gỗ thì
\choice
{lực của búa tác dụng vào đinh lớn hơn lực đinh tác dụng vào búa}
{\True lực của búa tác dụng vào đinh về độ lớn bằng lực của đinh tác dụng vào búa}
{lực của búa tác dụng vào đinh nhỏ hơn lực đinh tác dụng vào búa}
{tùy thuộc đinh di chuyển nhiều hay ít mà lực do đinh tác dụng vào búa lớn hơn hay nhỏ hơn lực do búa tác dụng vào đinh}
\loigiai{
Theo định luật III Newton: Lực của búa tác dụng vào đinh về độ lớn bằng lực của đinh tác dụng vào búa

}
\haidong{8} \end{ex}
\begin{ex}%book
Để xách một túi đựng thức ăn, một người tác dụng vào túi một lực bằng $40\ N$ hướng lên trên. Độ lớn của phản lực và hướng của phản lực (theo định luật 3) đạt giá trị bằng bao nhiêu và được xác định như thế nào?
\choice
{\True $40\ N$, hướng xuống dưới (ngược với chiều người tác dụng)}
{$50\ N$, hướng lên trên (ngược với chiều người tác dụng)}
{$40\ N$, hướng lên trên (cùng với chiều người tác dụng)}
{$50\ N$, hướng xuống dưới (ngược với chiều người tác dụng)}
\loigiai{
Theo định luật III Newton

$\Rightarrow F_{21}=F_{12}=40\ N$

 Độ lớn của phản lực là $40\ N$

}
\haidong{8} \end{ex}
\begin{ex}%book
Một người có trọng lượng $500\ N$ đứng trên mặt đất. Lực mà mặt đất tác dụng lên người đó có độ lớn
\choice
{\True bằng $500\ N$}
{lớn hơn $500\ N$}
{nhỏ hơn $500\ N$}
{bằng $250\ N$}
\loigiai{
\begin{center}
\includegraphics[width=4cm]{images/2023_11_19_566303217d153e753fb8g-08}
\end{center}	
}
\haidong{8} \end{ex}
\loai{Tổng hợp}
\begin{ex}%book%[1P2B4-1][Loai1]
Quả bóng đập vào tường bật ngược lại được là do nguyên nhân nào dưới đây?
\choice
{Lực do bóng tác dụng lên tường}
{\True Phản lực do tường tác dụng lên bóng}
{Trọng lực của bóng}
{Quán tính của bóng}
\haidong{8} \end{ex}

\begin{ex}%book%[1P2B4-1][Loai1]
Khi một con ngựa kéo xe, lực tác dụng vào con ngựa làm cho nó chuyển động về phía trước là
\choice
{lực mà con ngựa tác dụng vào xe}
{lực mà xe tác dụng vào ngựa}
{lực mà ngựa tác dụng vào đất}
{\True lực mà đất tác dụng vào ngựa}
\haidong{8} \end{ex}
\begin{ex}%book
Khi một người đẩy một thùng hàng chuyển động, lực tác dụng vào người làm người đó chuyển động về phía trước là
\choice
{lực người tác dụng vào thùng hàng}
{lực mà thùng hàng tác dụng vào người}
{lực người tác dụng vào mặt đất}
{\True lực mặt đất tác dụng vào người}
\loigiai{
Chọn D
}
\haidong{8} \end{ex}

%\newpage


\begin{ex}%book
Hình bên vẽ các lực tác dụng (cùng tỉ lệ) lên một chiếc xe đang chuyển động trên sàn ngang theo chiều dương. Phát biểu nào sau đây đúng?
\immini{\choice
{\True $\vec{N}$ là phản lực của sàn tác dụng lên xe}
{Xe có thể đang chuyển động chậm dần}
{$\vec{F}_{k}$ và $\vec{F}_C$ không có phản lực}
{Chỉ có lực $\vec{F}_{k}$ gây ra gia tốc cho xe}}{
\begin{tikzpicture}[scale=1,thick,>=stealth',draw=Mapcolor]
\tkzDefPoints{0/0/o, -0.35/-0.25/a, 0.35/-0.25/b} \tkzDrawPoints[size=3,fill=white](o)
\draw[Mapcolor] (-1.1,-0.4) -- (1.1,-0.4);
\fill[pattern=north east lines] (-1.1,-0.4)rectangle(1.1,-0.5);
\node  at  (o)  [shape=rectangle,draw=red,fill=red!40!white,minimum height=.45cm,minimum width=1.25cm, rounded corners=1pt]  {};
\node [draw, thick, shape=circle,draw=red,fill=white, minimum width=.1cm, minimum height=.1cm, anchor=center] at (a) {};
\node [draw, thick, shape=circle,draw=red,fill=white, minimum width=.1cm, minimum height=.1cm, anchor=center] at (b) {};
\draw[<->,Mapcolor,very thick] (-1,0) -- (1.4,0);
\draw [<->,Mapcolor,very thick](0,1) -- (0,-1);
\node at (0.2,1.1) {$\vec N$};
\node at (0.2,-0.9) {$\vec P$};
\node at (-1.3,0) {$\vec F$};
\node at (1.7,0) {$\vec F_k$};
\draw[->,Mapcolor] (0.8,0.5) -- (1.6,0.5);
\node at (2,0.5) {(+)};
\end{tikzpicture}}
\loigiai{
Chọn A
}
\haidong{8} \end{ex}
\begin{ex}%book
Hình bên vẽ các lực tác dụng lên một chiếc xe đang chuyển động với vận tốc v trên đường ngang. Phát biểu nào sau đây đúng?
\immini{\choice
{$\vec{P}$ không có phản lực}
{$\vec{F}$ không có phản lực}
{$\vec{N}$ và $\vec{P}$ là hai lực trực đối}
{\True $\vec{F}$ là lực cản chuyển động của xe}}{
\begin{tikzpicture}[scale=1,thick,>=stealth']
\tkzDefPoints{0/0/o, -0.35/-0.25/a, 0.35/-0.25/b} \tkzDrawPoints[size=3,fill=white](o)
\draw[->,Mapcolor,very thick] (-1.1,-0.4) -- (1.6,-0.4);


\node  at  (o)  [shape=rectangle,draw=red,fill=red!40!white,minimum height=.45cm,minimum width=1.25cm, rounded corners=1pt]  {};
\node [draw, thick, shape=circle,draw=red,fill=white, minimum width=.1cm, minimum height=.1cm, anchor=center] at (a) {};
\node [draw, thick, shape=circle,draw=red,fill=white, minimum width=.1cm, minimum height=.1cm, anchor=center] at (b) {};
\draw[<->,Mapcolor,very thick] (-1,0) -- (1,0);
\draw [<->,Mapcolor,very thick](0,1) -- (0,-1);
\node at (0.2,1.1) {$\vec N$};
\node at (0.2,-0.9) {$\vec P$};
\node at (-1.3,0) {$\vec F$};
\node at (1,0.2) {$\vec v$};
\end{tikzpicture}}

\loigiai{

F ngược chiều chuyển động nên là lực cản chuyển động của xe

}
\haidong{8} \end{ex}

%mac
%\loai{Va chạm vuông góc với tường}






\subsubsection*{PHẦN 2. Thí sinh trả lời từ câu 1 đến câu 4. Trong mỗi ý a), b), c), d) ở mỗi câu, thí sinh chọn đúng hoặc sai.}
\setcounter{ex}{0}
\begin{ex}%book
Xét tương tác giữa hai vật.
\choiceTF[t]
{Gia tốc mà hai vật thu được luôn ngược chiều nhau và có độ lớn tỉ lệ thuận với khối lượng của chúng}
{\True Hai lực trực đối đặt vào hai vật khác nhau nên không cân bằng nhau}
{\True Các lực tương tác giữa hai vật là hai lực trực đối}
{\True Lực và phản lực có độ lớn bằng nhau}
\loigiai{
}
\haidong{8} \end{ex}
\begin{ex}%book
Xét các đặc điểm của lực và phản lực.
\choiceTF[t]
{\True Lực và phản lực luôn xuất hiện thành từng cặp (xuất hiện hoặc mất đi đồng thời)}
{\True Lực và phản lực cùng tác dụng theo một đường thẳng, cùng độ lớn nhưng ngược chiều (hay còn gọi là hai lực trực đối)}
{\True Lực và phản lực không cân bằng nhau (vì chúng đặt vào hai vật khác nhau)}
{Cặp lực và phản lực là hai lực khác loại}
\loigiai{}
\haidong{8} \end{ex}
\begin{ex}%book
Hình bên vẽ các lực tác dụng lên một chiếc xe đang chuyển động với vận tốc $\vec{v}$ trên đường ngang.
\immini{\choiceTF[t]
{$\vec{N}$ và $\vec{P}$ là lực và phản lực}
{\True Xe đang chuyển động chậm dần}
{\True $\vec{N}$ và $\vec{P}$ là hai lực cân bằng}
{\True Chỉ có lực $\vec{F}$ gây ra gia tốc cho xe}}{
\begin{tikzpicture}[scale=1,thick,>=stealth',draw=Mapcolor]
\tkzDefPoints{0/0/o, -0.35/-0.25/a, 0.35/-0.25/b} \tkzDrawPoints[size=3,fill=white](o)
\draw[->,Mapcolor,very thick] (-1.1,-0.4) -- (1.6,-0.4);


\node  at  (o)  [shape=rectangle,draw=red,fill=red!40!white,minimum height=.45cm,minimum width=1.25cm, rounded corners=1pt]  {};
\node [draw, thick, shape=circle,draw=red,fill=white, minimum width=.1cm, minimum height=.1cm, anchor=center] at (a) {};
\node [draw, thick, shape=circle,draw=red,fill=white, minimum width=.1cm, minimum height=.1cm, anchor=center] at (b) {};
\draw[<->,Mapcolor,very thick] (-1,0) -- (1,0);
\draw [<->,Mapcolor,very thick](0,1) -- (0,-1);
\node at (0.2,1.1) {$\vec N$};
\node at (0.2,-0.9) {$\vec P$};
\node at (-1.3,0) {$\vec F$};
\node at (1,0.2) {$\vec v$};
\end{tikzpicture}}
\loigiai{
Chọn A

$N$ là phản lực của mặt sàn tác dụng lên xe

\begin{center}
\includegraphics[width=4cm]{images/2023_11_19_566303217d153e753fb8g-06(1)}
\end{center}

}
\haidong{8} %<MyLT1>
\end{ex}



















\begin{ex}%book
Một vật có khối lượng $M=33\ kg$ được đẩy bằng một thanh sắt có khối lượng $m=3\ kg$. Vật chuyển động (từ trạng thái đứng yên) một đoạn 80 cm trong thời gian $1,6\ s$ với gia tốc không đổi.
\choiceTF[t]
{\True Khi không có ma sát, hệ gồm thanh sắt và vật chuyển động nhanh dần đều}
{\True Có một cặp lực và phản lực theo phương ngang}
{Tay sẽ phải tác dụng lên thanh sắt một lực bằng $19,3\ N$}
{Hợp lực tác dụng lên thanh sắt bằng $43,125\ N$}
\loigiai{\begin{itemchoice}
\itemch Đúng.
\itemch \immini{b) Đúng. Cặp thứ nhất: $\vec{F}_{1}$ (lực do tay tác dụng lên thanh sắt) và $\vec{F}_{1}^{\prime}$ (lực do thanh sắt tác dụng lên tay)\\
Cặp thứ hai: $\vec{F}_{2}$ (lực do thanh sắt tắc dụng lên vật) và $\vec{F}_{2}^{\prime}$ (lực do vật tác dụng lên thanh sắt)}{\begin{tikzpicture}[scale=1,thick,>=stealth']
\tkzDefPoints{0/1/o}
\node  at  (-1.5,0.8)  [shape=rectangle,draw=black,fill=orange,minimum height=1cm,minimum width=1.8cm, rounded corners=1pt]  {$1$};
\draw[snake=coil,segment length=3pt,segment amplitude=5.5pt,line before snake=3pt,line after snake=3pt] (-0.6,0.8) -- (0.6,0.8);
\draw[very thick](-0.6,0.73) -- (0.6,0.73);
\node  at  (1.5,0.8)  [shape=rectangle,draw=black,fill=green,minimum height=1cm,minimum width=1.8cm, rounded corners=1pt]  {$2$};
\draw(-2.5,0)--(2.5,0);
\draw[fill=white](-2,0.2)circle(0.2cm);
\draw[fill=white](-2,0.2)circle(0.06cm);
\draw[fill=white](-1,0.2)circle(0.2cm);
\draw[fill=white](-1,0.2)circle(0.06cm);
\draw[fill=white](2,0.2)circle(0.2cm);
\draw[fill=white](2,0.2)circle(0.06cm);
\draw[fill=white](1,0.2)circle(0.2cm);
\draw[fill=white](1,0.2)circle(0.06cm);
\fill[pattern=north east lines](-2.5,0)--(2.5,0)--++(-90:0.1)--++(180:5);
\end{tikzpicture}}
\itemch Sai. Gia tốc của thanh sắt và vật $S=v_{0} t+\dfrac{1}{2} a t^{2} \Rightarrow a=\dfrac{2 S}{t^{2}}=\dfrac{2 . 0,8}{1,6^{2}}=0,625\ m/s^{2}$.\\
Áp dụng định luật II Newton cho hệ gồm thanh sắt và vật: $F_{1}=(M+m) a=(33+3) . 0,625=22,5\ N$.
\itemch Sai. Áp dụng định luật II Newton cho vật khối lượng $M: F_{2}=M a=33 . 0,625=20,625\ N$.\\
Theo định luật III Newton thì $F_{2}^{\prime}=F_{2}=20,625\ N$.\\
Hai lực $\vec{F}_{1}$ và $\vec{F}_{2}{ }^{\prime}$ tác dụng lên thanh sắt cùng phương, ngược chiều nhau, độ lớn hợp lực của hai lực này là: $F=F_{1}-F_{2}^{\prime}=22,5-20,625=1,875\ N$.
\end{itemchoice}
}
\haidong{7} \end{ex}
\subsubsection*{PHẦN 3. Thí sinh trả lời từ câu 1 đến câu 6.}
\setcounter{ex}{0} 
\begin{ex}%book%[1P2K4-2][Loai1]
Một quả bóng khối lượng $m = 200\ g$ bay với vận tốc $v_1 = 20\ m/s$ thì đập vuông góc vào một bức tường rồi bật trở lại theo phương cũ với vận tốc $v_2= 10\ m/s$. Khoảng thời gian va chạm vào tường là $\Delta t = 0,05\ s$. Lực của tường tác dụng lên quả bóng có độ lớn bằng bao nhiêu niutơn?
\shortans[oly]{120}
\loigiai{
\begin{itemize}
\item Gia tốc của quả bóng là: $a = \dfrac{v_2 - v_1}{t} = \dfrac{10-(-20)}{0,05} = 600\ m/s^2$.
\item Lực do tường tác dụng lên quả bóng là: $F=m.a=0,2.600=120\ N$.
\end{itemize}
}
\haidong{7} \end{ex}
\begin{ex}%book%[1P2K4-2][Loai1]
Quả bóng khối lượng $500\ g$ bay với vận tốc $72\ km/h$ đến đập vuông góc vào một bức tường rồi bật trở ra theo phương cũ với vận tốc $54\ km/h$. Thời gian va chạm là $0,05\ s$. Lực của bóng tác dụng lên tường có độ lớn bằng bao nhiêu niutơn?
\shortans[oly]{350}
\haidong{12}\loigiai{
\begin{itemize}
\item Gia tốc của quả bóng là: $a = \dfrac{v_2-v_1}{t} = \dfrac{15-(-20)}{0,05} = 700\ m/s^2$.
\item Lực do tường tác dụng lên quả bóng là: $F=m.a=0,5.700=350\ N$.
\end{itemize}
}
\haidong{7} \end{ex}
\begin{ex}%Tailieuchuan
Quả bóng khối lượng $200\ g$ bay với vận tốc $90\ km/h$ đến đập vuông góc vào một bức tường rồi bật trở lại theo phương cũ với vận tốc $54\ km/h$. Thời gian va chạm là $0,05\ s$. Lực do tường tác dụng lên quả bóng có độ lớn bao nhiêu niutơn?
\shortans[oly]{160}
\loigiai{}
\end{ex}
\begin{ex}%book
Một quả bóng có khối lượng $0,2\ kg$ bay với tốc độ $25\ m/s$ đến đập vuông góc tường rồi bật ngược trở lại theo phương cũ với tốc độ $15\ m/s$. Khoảng thời gian va chạm $0,1\ s$. Coi lực này là không đổi trong suốt thời gian tác dụng. Lực của tường tác dụng lên quả bóng có độ lớn bằng bao nhiêu niutơn?
\shortans[oly]{80}
\haidong{10}
\loigiai{
Chọn chiều dương là chiều chuyển động của quả bóng sau khi va chạm.

Lực của tường tác dụng lên quả bóng:

$F_{21}=m_1 a_1=m_1 \dfrac{\Delta v_1}{\Delta t}=0,2 \cdot \dfrac{15-(-25)}{0,1}=80(~N)$

}
\haidong{7} \end{ex}
\begin{ex}%book%[1P2K4-2][Loai1]
Quả bóng khối lượng $500\ g$ bay với vận tốc $72\ km/h$ đến đập vào tường với góc tới $60^\circ$ và bật trở lại với độ lớn vận tốc không đổi. Biết va chạm của bóng với tường tuân theo định luật phản xạ gương (góc phản xạ bằng góc tới), thời gian va chạm là $0,05\ s$. Lực do bóng tác dụng lên tường có độ lớn bằng bao nhiêu niutơn?
\shortans[oly]{200}
\loigiai{
\begin{itemize}
\item Đổi $72\ km/h = 20\ m/s$.
\item Vectơ gia tốc của quả bóng bàn $\vec a = \dfrac{\vec v'-\vec v}{\Delta t} = \dfrac{\Delta\vec{v}}{\Delta t}$.
\item Từ hình vẽ $\Rightarrow \Delta v = 2v\cos 60^\circ = 2.20.\cos60^\circ = 20 \ m/s$. 
\item Lực do tường tác dụng lên bóng có độ lớn là: $\vec F = m\vec a = m\dfrac{|\Delta\vec{v}|}{\Delta t} = 0,5.\dfrac{20}{0,05}=200$ N.
\end{itemize}
}
\haidong{7} \end{ex}
\begin{ex}%Tailieuchuan
Một quả bóng chày có khối lượng $300\ g$ bay với vận tốc $72\ km/h$ đến đập vuông góc với tường và bật ngược trở lại theo phương cũ với vận tốc $54\ km/s$. Thời gian va chạm là $0,04\ s$. Lực do tường tác dụng vào quả bóng có độ lớn bao nhiêu niutơn (làm tròn kết quả đến chữ số hàng đơn vị)?
\shortans[oly]{263}
\loigiai{}
\end{ex}
